% !TEX program = pdflatex
% Amaris Consulting — Data Engineer / Data Ingestion Factory · Tunis · FR
% Points forts : Python, Pandas-minded data flows, APIs REST, MongoDB/SQL, CI/CD, Docker, Agile, FR+EN, Tunis
% Écarts honnêtes : 4–6 ans exigés (profil junior) ; pas de PySpark/Kafka/Databricks en prod — ne pas inventer
\documentclass[10pt,a4paper]{article}

\usepackage[french]{babel}
\usepackage[utf8]{inputenc}
\usepackage[T1]{fontenc}
\usepackage{lmodern}
\usepackage[margin=1.05cm]{geometry}
\usepackage{xcolor}
\usepackage{titlesec}
\usepackage{enumitem}
\usepackage{hyperref}
\usepackage{fontawesome5}
\usepackage{microtype}

\setlength{\parindent}{0pt}
\setlength{\parskip}{1pt}
\pagestyle{empty}
\setlist[itemize]{leftmargin=0.85em, itemsep=0.5pt, topsep=0pt, parsep=0pt, partopsep=0pt}

\definecolor{primary}{HTML}{1A365D}
\definecolor{accent}{HTML}{2B6CB0}
\definecolor{muted}{HTML}{4A5568}
\definecolor{rulecolor}{HTML}{E2E8F0}

\newcommand{\TargetTitle}{Data Engineer junior \textbar{} Python · Pipelines · APIs · CI/CD}
\newcommand{\TargetKeywords}{Python · Pandas · API REST · MongoDB · PostgreSQL · Docker · Git · CI/CD · Agile · Ingestion · Qualité code · Anglais · Français}

\hypersetup{
  colorlinks=true,
  linkcolor=accent,
  urlcolor=accent,
  pdfauthor={Hamouda Chekir},
  pdftitle={CV — Hamouda Chekir — Amaris Data Engineer}
}

\titleformat{\section}
  {\color{primary}\normalsize\bfseries}
  {}
  {0pt}
  {}
  [\vspace{0.5pt}{\color{accent}\titlerule[1.4pt]}\vspace{3pt}]
\titlespacing*{\section}{0pt}{7pt}{2pt}

\newcommand{\job}[3]{%
  \textbf{\color{primary}#1}\hfill{\color{muted}\small #2}\\[-1pt]
  {\color{accent}\textit{#3}}\\[1pt]
  \begin{itemize}
}

\newcommand{\proj}[3]{%
  \textbf{\color{primary}#1}\hfill{\color{muted}\small #2}\\[-1pt]
  {\color{accent}\textit{#3}}\\[0.5pt]
  \begin{itemize}
}

\newcommand{\edu}[3]{%
  \textbf{\color{primary}#1}\hfill{\color{muted}\small #2}\\[-1pt]
  {\color{accent}\textit{#3}}\\[2pt]
}

\newcommand{\contactsep}{\hspace{4pt}{\color{rulecolor}|}\hspace{4pt}}

\begin{document}

{\fontsize{20}{24}\selectfont\bfseries\color{primary}HAMOUDA CHEKIR}\\[1pt]
{\normalsize\color{accent}\TargetTitle}\\[3pt]
{\footnotesize\color{muted}%
  \faIcon{envelope}\hspace{2pt}\href{mailto:hamoudachkir@yahoo.fr}{hamoudachkir@yahoo.fr}%
  \contactsep
  \faIcon{phone}\hspace{2pt}+216\,55\,913\,832%
  \contactsep
  \faIcon{map-marker-alt}\hspace{2pt}Tunis, Tunisie%
  \contactsep
  \faIcon{linkedin}\hspace{2pt}\href{https://www.linkedin.com/in/hamouda-chekir-5053572b4}{LinkedIn}%
  \contactsep
  \faIcon{github}\hspace{2pt}\href{https://github.com/hamoudachekir}{GitHub}%
}\\[4pt]

\section{PROFIL}
Ingénieur junior (ESPRIT — TWIN) orienté \textbf{Python}, pipelines et APIs : ingestion, transformation et persistance (MongoDB, PostgreSQL) avec qualité, tests et maintenabilité.
Chez Talan : services data/IA sur NextHire (\textbf{16 microservices} Dockerisés) — parsing CV, intégrations REST, CI/CD et troubleshooting de parcours critiques.
Français et anglais courants ; Agile, code review, documentation technique. Motivé à monter sur \textbf{PySpark, Kafka/Confluent et Databricks} en Data Ingestion Factory.

\section{EXPÉRIENCE}

\job{Ingénieur logiciel / data — Projet de fin d'études}{Fév.\ 2026 -- Août 2026}{Talan \textbar{} Tunis}
  \item Conçu et maintenu des pipelines Python (FastAPI/Flask) d'ingestion/traitement (parsing CV spaCy/OCR, scoring) vers MongoDB, exposés via APIs REST.
  \item Industrialisé le delivery : Docker, GitLab CI/CD, revues de code, hotfix/troubleshooting sur parcours temps réel.
  \item Collaboré en Agile avec plusieurs équipes ; documenté contrats d'API pour faciliter onboarding et run.
\end{itemize}

\job{Stagiaire développeur Full-Stack}{Juin 2025 -- Août 2025}{SmartConseil — Dafe Management \textbar{} Tunis}
  \item Intégré des flux de données (CRUD, exports PDF/Excel) sur PostgreSQL/MongoDB ; validation via tests, SonarQube, GitLab CI/CD et Docker.
\end{itemize}

\job{Stagiaire développeur web}{Juin 2024 -- Août 2024}{Big Deal (Club Privilèges) \textbar{} Tunis}
  \item Livré un module métier avec validations API (Postman) et Git avant mise en production.
\end{itemize}

\section{PROJETS SÉLECTIONNÉS}
{\footnotesize
\proj{Moteur AML — ingestion \& traçabilité}{2025}{Python · règles métier · RAG · MCP}
  \item Pipeline d'analyse (règles + agents LLM), gateway sécurisée et traçabilité — esprit qualité / audit proche d'un run data.
\end{itemize}
\proj{TalentMatch — APIs data \& SaaS}{2025}{Python/FastAPI · scraping · APIs externes}
  \item Orchestré connecteurs API/SaaS et traitements Python (authenticité, pricing, fiches) — structuration et livrables fiables.
\end{itemize}
}

\section{FORMATION}
\edu{Diplôme d'ingénieur — TWIN (Technologies du Web, de l'Internet et des Réseaux)}{2023 -- 2026}{ESPRIT — École Supérieure Privée d'Ingénierie et de Technologies \textbar{} Tunis}
\edu{Licence en informatique de gestion — E-Business}{2020 -- 2023}{ESEN — École Supérieure d'Économie Numérique \textbar{} Manouba}

\section{COMPÉTENCES}
{\small
\textbf{\color{primary}Cœur data / ingestion :} Python, Pandas (pratique), APIs REST, MongoDB, PostgreSQL/MySQL, fichiers \& exports, Docker\\[1pt]
\textbf{\color{primary}Industrialisation :} Git, GitLab CI/CD (proche GitHub Actions), tests / revues, packaging services, Linux/Bash\\[1pt]
\textbf{\color{primary}Run \& collab :} troubleshooting, Agile, documentation technique FR/EN, équipes multi-parties\\[1pt]
\textbf{\color{primary}Montée en compétence :} PySpark, Kafka/Confluent, Databricks, observabilité FinOps — prêt à ramper en Data Factory
}

\vspace{1pt}
\begin{minipage}[t]{0.70\textwidth}
\section{CERTIFICATIONS}
{\small
\textbf{\color{primary}AWS Academy Graduate — Cloud Foundations} — AWS Academy (mai 2025)\\[1pt]
\textbf{\color{primary}CCNA — Switching, Routing \& Wireless Essentials} — Cisco (sept.\ 2024)
}
\end{minipage}%
\hfill
\begin{minipage}[t]{0.26\textwidth}
\section{LANGUES}
{\small\textbf{Arabe} — Langue maternelle}\\[1pt]
{\small\textbf{Français} — Courant}\\[1pt]
{\small\textbf{Anglais} — Courant}
\end{minipage}

\end{document}
